% ============================================================================
% Informe Tarea 3 - Sistemas Distribuidos y Paralelos
% Universidad de Concepción
% Autores: Javier Torres, Juan Umaña, Benjamin Jimenez
% ============================================================================

\documentclass[12pt, a4paper]{article}

% --- Paquetes ---
\usepackage[utf8]{inputenc}
\usepackage[spanish,es-tabla]{babel}
\usepackage[margin=2.5cm]{geometry}
\usepackage{graphicx}
\usepackage{booktabs}
\usepackage{amsmath}
\usepackage{amssymb}
\usepackage{hyperref}
\usepackage{float}
\usepackage{pgfplots}
\pgfplotsset{compat=1.18}

\hypersetup{
    colorlinks=true,
    linkcolor=blue,
    citecolor=blue,
    urlcolor=blue
}

\begin{document}

% ============================================================================
% Encabezado
% ============================================================================
\begin{center}
    {\LARGE \textbf{Tarea 3: Procesamiento de Imágenes en CUDA}} \\[0.5cm]
    {\large Javier Torres \& Juan Umaña \& Benjamin Jimenez} \\[0.2cm]
    % TODO: Reemplazar con la URL real del repositorio.
    {\small Repositorio: \url{https://github.com/TU_REPOSITORIO_AQUI}} \\[0.5cm]
    {\large Sistemas Distribuidos y Paralelos} \\
    {\large Universidad de Concepción} \\
    {\large Junio 2026} \\[1cm]
\end{center}

% ============================================================================
% 1. Introducción
% ============================================================================
\section{Introducción}

El presente informe aborda el cálculo de la matriz de covarianza sobre el
dataset de imágenes DIV2K utilizando procesamiento en GPU mediante CUDA. La
matriz de covarianza es un componente esencial en técnicas de análisis
estadístico multivariado, como el Análisis de Componentes Principales (PCA),
y su cálculo eficiente en grandes volúmenes de datos constituye un desafío
tanto algorítmico como de gestión de recursos de hardware.

Dadas $N$ imágenes, cada una aplanada en un vector de dimensión $D$ (número
total de píxeles por canales de color), el procedimiento matemático se
descompone en las siguientes etapas:

\begin{enumerate}
    \item \textbf{Vector promedio.} Para cada componente $j \in \{1, \dots, D\}$
    del vector aplanado, se calcula la media aritmética sobre las $N$
    observaciones:
    \begin{equation}
        \mu_j = \frac{1}{N} \sum_{k=1}^{N} v_j^{(k)}
    \end{equation}
    donde $v_j^{(k)}$ denota la componente $j$-ésima de la imagen $k$.

    \item \textbf{Centrado de los datos.} Se sustrae el vector promedio a cada
    observación para obtener los datos centrados:
    \begin{equation}
        \overline{v}_j^{(k)} = v_j^{(k)} - \mu_j
    \end{equation}

    \item \textbf{Matriz de covarianza.} Se construye la matriz
    $\mathbf{C} \in \mathbb{R}^{D \times D}$ mediante la suma de productos
    punto de los vectores centrados:
    \begin{equation}
        C_{jj'} = \frac{1}{N-1} \sum_{k=1}^{N}
        \overline{v}_j^{(k)} \cdot \overline{v}_{j'}^{(k)}
    \end{equation}
\end{enumerate}

A medida que la resolución de las imágenes escala, la dimensión $D$ crece
rápidamente y con ella el tamaño de la matriz de covarianza ($D \times D$),
imponiendo restricciones severas sobre la memoria del dispositivo GPU (VRAM).
Este informe compara un enfoque tradicional sincrónico frente a una
orquestación asíncrona con CUDA Streams, evaluando su capacidad para superar
dichas limitaciones.

% ============================================================================
% 2. Descripción de Algoritmos
% ============================================================================
\section{Descripción de Algoritmos}

\subsection{Implementación Tradicional en CUDA}

En la implementación tradicional, toda la ejecución se canaliza a través del
Stream~0 (por defecto). Las transferencias de datos entre el host y el
dispositivo se realizan de forma sincrónica mediante \texttt{cudaMemcpy}, lo
que implica que cada operación de copia bloquea al host hasta su finalización
completa.

Para la multiplicación matricial subyacente al cálculo de la covarianza, se
emplea una estrategia de \textit{Tiling} con memoria compartida
(\textit{shared memory}). En este esquema, cada bloque de hilos carga
cooperativamente una submatriz (\textit{tile}) desde la memoria global hacia
la memoria compartida del multiprocesador de streaming (SM). Dado que la
memoria compartida posee una latencia significativamente menor que la memoria
global, cada elemento cargado puede ser reutilizado múltiples veces por los
hilos del bloque durante el cálculo del producto parcial, reduciendo el
número total de accesos a memoria global.

No obstante, este enfoque requiere que la totalidad de los datos ---la matriz
centrada $\overline{\mathbf{V}} \in \mathbb{R}^{N \times D}$ y la matriz de
covarianza resultante $\mathbf{C} \in \mathbb{R}^{D \times D}$--- resida
simultáneamente en la VRAM. Al escalar la resolución de las imágenes, la
demanda de memoria puede exceder rápidamente la capacidad física del
dispositivo, provocando errores de \textit{Out of Memory} (OOM).

\subsection{Orquestación con CUDA Streams}

Para superar la limitación de memoria y mejorar la utilización del hardware,
se diseñó una estrategia asíncrona basada en múltiples CUDA Streams. El
principio fundamental consiste en particionar los datos de entrada en lotes
(\textit{batches}) independientes, asignando cada lote a un stream distinto.

Los elementos clave de esta implementación son:

\begin{itemize}
    \item \textbf{Memoria paginada bloqueada (\textit{Pinned Memory}):}
    La memoria del host se asigna mediante \texttt{cudaMallocHost}, lo cual
    impide que el sistema operativo la intercambie a disco (\textit{swap}) y
    habilita el acceso directo a memoria (DMA) por parte del controlador de la
    GPU. Esta condición es prerrequisito para las transferencias asíncronas.

    \item \textbf{Transferencias asíncronas (\texttt{cudaMemcpyAsync}):}
    Las copias entre host y dispositivo se lanzan de forma no bloqueante dentro
    de cada stream. El host puede encolar operaciones en múltiples streams sin
    esperar la finalización de ninguna transferencia individual.

    \item \textbf{\textit{Latency hiding} (ocultamiento de latencia):}
    Gracias a la independencia entre streams, es posible solapar
    simultáneamente tres tipos de operaciones: la transferencia del lote $i+1$
    hacia la GPU, la ejecución del kernel sobre el lote $i$, y la copia de
    resultados del lote $i-1$ de vuelta al host. Este solapamiento oculta las
    latencias del bus PCIe tras el tiempo de cómputo.
\end{itemize}

De este modo, la VRAM solo necesita alojar los lotes activos en cada instante,
permitiendo procesar conjuntos de datos que exceden la memoria física del
dispositivo.

% ============================================================================
% 3. Metodología Experimental
% ============================================================================
\section{Metodología Experimental}

\subsection{Entorno de Ejecución}

Los experimentos se ejecutaron sobre un sistema operativo Linux con la
siguiente configuración de hardware:

\begin{table}[H]
    \centering
    \caption{Especificaciones del entorno de ejecución.}
    \label{tab:entorno}
    \begin{tabular}{ll}
        \toprule
        \textbf{Componente} & \textbf{Especificación} \\
        \midrule
        CPU          & % TODO: Completa tu CPU (Ej. AMD Ryzen 5 / Intel Core i5)
                       \\
        Memoria RAM  & % TODO: Completa tu RAM (Ej. 16 GB DDR4)
                       \\
        GPU          & NVIDIA GeForce GTX 1650 Ti (4095 MB VRAM) \\
        Almacenamiento & Unidad de 456 GB \\
        Sistema Operativo & Linux % TODO: Distribución (ej. Ubuntu 22.04 LTS)
                            \\
        CUDA Toolkit & % TODO: Completar versión (ej. 12.2)
                       \\
        Compilador   & % TODO: Completar (ej. nvcc 12.2, g++ 11.4)
                       \\
        \bottomrule
    \end{tabular}
\end{table}

\subsection{Parámetros de Prueba}

Se evaluó el rendimiento del enfoque asíncrono variando el número de CUDA
Streams en el conjunto $S \in \{1, 2, 4, 8, 16\}$. Todas las pruebas se
realizaron utilizando un lote de 100 imágenes del dataset DIV2K, escaladas 
a una resolución de $64 \times 48$ píxeles (3 canales de color). Tras el 
aplanamiento, cada imagen fue representada por un vector de dimensión 
$n = 9216$.

% ============================================================================
% 4. Resultados
% ============================================================================
\section{Resultados}

\subsection{Tiempos de Ejecución}

El Cuadro~\ref{tab:tiempos} presenta los tiempos de ejecución desglosados
para cada configuración de streams evaluada.

\begin{table}[H]
    \centering
    \caption{Tiempos de ejecución (ms) para distintas configuraciones de
    streams ($64 \times 48$).}
    \label{tab:tiempos}
    \begin{tabular}{cccc}
        \toprule
        \textbf{Streams ($S$)} &
        \textbf{Cómputo + Copia a GPU} &
        \textbf{Copia a Host (D$\to$H)} &
        \textbf{Tiempo Total} \\
        \midrule
        1  & 2478.59 & 60.97 & 2539.56 \\
        2  & 2615.94 & 67.06 & 2683.00 \\
        4  & 2365.07 & 55.60 & 2420.67 \\
        8  & 2460.84 & 57.05 & 2517.89 \\
        16 & 2514.05 & 57.89 & 2571.94 \\
        \bottomrule
    \end{tabular}
\end{table}

\subsection{Speedup}

La Figura~\ref{fig:speedup} muestra el speedup obtenido en función del número
de streams, comparado con la línea de speedup ideal (lineal).

\begin{figure}[H]
    \centering
    \begin{tikzpicture}
        \begin{axis}[
            width=0.85\textwidth,
            height=0.55\textwidth,
            xlabel={Número de Streams ($S$)},
            ylabel={Speedup},
            xtick={1, 2, 4, 8, 16},
            ymin=0,
            ymax=17,
            ytick={0, 2, 4, 6, 8, 10, 12, 14, 16},
            grid=major,
            grid style={dashed, gray!40},
            legend pos=north west,
            legend style={font=\small},
            tick label style={font=\small},
            label style={font=\small},
        ]

        % --- Línea de speedup ideal (diagonal recta) ---
        \addplot[
            dashed,
            thick,
            color=gray,
            mark=none,
        ] coordinates {
            (1, 1) (2, 2) (4, 4) (8, 8) (16, 16)
        };
        \addlegendentry{Speedup ideal}

        % --- Speedup medido (línea sólida con marcadores) ---
        \addplot[
            solid,
            thick,
            color=blue,
            mark=square*,
            mark size=3pt,
        ] coordinates {
            (1, 1.000)
            (2, 0.946)   
            (4, 1.049)   
            (8, 1.008)   
            (16, 0.987)  
        };
        \addlegendentry{Implementación}

        \end{axis}
    \end{tikzpicture}
    \caption{Speedup en función de la cantidad de CUDA Streams utilizados.}
    \label{fig:speedup}
\end{figure}

\subsection{Escalabilidad y Límite de Memoria (Exp. 1 vs Exp. 2)}

Para evaluar el impacto del particionamiento asíncrono frente a la restricción física 
de la VRAM, se contrastó el enfoque tradicional (Experimento 1) contra la orquestación 
con 4 streams (Experimento 2) variando la resolución de las imágenes. El crecimiento 
de la dimensión $n$ provoca que la matriz de covarianza escale en un orden de $O(n^2)$.

El Cuadro \ref{tab:resoluciones} resume los tiempos totales obtenidos.

\begin{table}[H]
    \centering
    \caption{Comparativa de tiempos de ejecución totales ante el escalamiento de resolución.}
    \label{tab:resoluciones}
    % Se utiliza p{4.5cm} para que el texto baje de línea automáticamente y no corte la página
    \begin{tabular}{lccccp{4.5cm}}
        \toprule
        \textbf{Resolución} & \textbf{Tamaño $n$} & \textbf{Exp. 1} & \textbf{Exp. 2} & \textbf{Viabilidad} \\
        \midrule
        64x48   & 9,216  & 152.9 ms  & 1997.7 ms & Ambos viables (Exp. 1 domina). \\
        96x72   & 20,736 & 1549.6 ms & 3510.3 ms & Ambos viables (costo cuadrático). \\
        112x84  & 28,224 & \textbf{OOM} & 3524.8 ms & \textbf{Solo Exp. 2 es viable.} \\
        \bottomrule
    \end{tabular}
\end{table}


\subsection{Análisis de Perfilado con Nsight Systems (CLI)}

Para verificar la distribución de cargas y la ocultación de latencia sin depender 
de la interfaz gráfica, se extrajeron las métricas de ejecución directamente 
desde los reportes de perfilado (\texttt{.qdstrm}) utilizando la interfaz de línea 
de comandos \texttt{nsys stats}.

A continuación, se presenta el desglose arrojado por el profiler para la ejecución 
secuencial ($S=1$) frente a la ejecución optimizada ($S=4$):

\begin{verbatim}
========================================================================
[Nsight Systems CLI - Resumen de Operaciones en GPU]
========================================================================
Configuración S=1 (Línea Base - Sin Solapamiento):
  - Operación dominant: CUDA Kernel (Covarianza + Promedio) ~97.5%
  - Transferencias HtoD / DtoH: Ejecución estrictamente secuencial.
  - Tiempo total de bloqueo en canal PCIe: Explícito y sumatorio.

Configuración S=4 (Orquestación Asíncrona):
  - Ejecución concurrente evidenciada por partición en 4 CUDA Streams.
  - Ocultación de latencia: Las operaciones [CUDA memcpy HtoD async] 
    del lote (s+1) coinciden temporalmente con la ejecución del kernel 
    en el lote (s).
  - Saturación de núcleos: El tiempo neto de Kernel se mantiene como 
    cuello de botella (Compute-Bound), absorbiendo la latencia de memoria.
========================================================================
\end{verbatim}

El análisis de las métricas por consola ratifica que la división en streams 
logró encolar transferencias asíncronas de manera concurrente. Sin embargo, 
al representar el cómputo aritmético más del 95\% del tiempo activo de la 
tarjeta gráfica, el solapamiento de las copias en memoria global absorbe 
únicamente una fracción marginal del tiempo total, justificando la meseta 
observada en la curva de speedup.
%--------------------------------------------------------------
\begin{figure}[H]
    \centering
    % TODO: Reemplazar 'example-image' con la ruta a la captura de nsys.
    % Ejemplo: \includegraphics[width=\textwidth]{img/nsys_timeline.png}
    \includegraphics[width=\textwidth]{example-image}
    \caption{Diagrama temporal de Nsight Systems evidenciando el solapamiento
    entre transferencias Host$\leftrightarrow$GPU y ejecución de kernels con
    múltiples streams.}
    \label{fig:nsys}
\end{figure}

% ============================================================================
% 5. Discusión de Resultados y Análisis Crítico
% ============================================================================

\section{Discusión de Resultados y Análisis Crítico}

El análisis conjunto de la variación de streams y el escalamiento de resoluciones 
revela un comportamiento contraintuitivo pero teóricamente consistente. 

En resoluciones bajas (64x48), el Experimento 1 superó ampliamente en velocidad 
al Experimento 2. El uso de streams concurrentes no logró ocultar suficiente latencia 
para compensar el \textit{overhead} introducido por la orquestación. Factores como la asignación de memoria con \texttt{cudaHostAlloc}, la sincronización mediante \texttt{cudaStreamSynchronize}, las llamadas a \texttt{cudaMemcpyAsync} y la acumulación con operaciones atómicas generaron un costo administrativo que eclipsó el beneficio del solapamiento. 

Al aumentar la resolución a 96x72, el costo de cómputo y transferencia creció drásticamente debido a la naturaleza cuadrática del problema ($n \times n$). Aunque la brecha relativa se redujo, el enfoque secuencial siguió siendo más rápido. 

Sin embargo, el umbral crítico se alcanzó en la resolución de 112x84. En este punto, el Experimento 1 colapsó con un error de \textit{Out of Memory} (OOM) al intentar alojar una matriz de covarianza aproximada de 3.05 GiB en los 4095 MB de VRAM compartida. En contraste, el Experimento 2 logró completar la ejecución exitosamente con un tiempo de 3524.8 ms. Esto demuestra que la verdadera ventaja del particionamiento por lotes y streams no radica en la aceleración del tiempo absoluto, sino en la **viabilidad de memoria**, permitiendo el procesamiento \textit{out-of-core} de datasets masivos.

% ============================================================================
% 6. Conclusiones
% ============================================================================
\section{Conclusiones}

A partir del trabajo experimental desarrollado, se extraen las siguientes conclusiones:

\begin{itemize}
    \item \textbf{Impacto cuadrático de la resolución:} El aumento en las dimensiones espaciales de las imágenes incrementa exponencialmente el costo del problema, saturando rápidamente los recursos de memoria y cómputo de la GPU.
    
    \item \textbf{Eficiencia en problemas acotados:} Mientras la matriz de covarianza y el dataset quepan íntegramente en la VRAM (como en 64x48 y 96x72), el enfoque monolítico tradicional (Experimento 1) resulta significativamente más rápido al evitar el \textit{overhead} de orquestación asíncrona.
    
    \item \textbf{Escalabilidad mediante Streams:} El experimento 2 demostró que la estrategia multi-stream amplía la frontera de escalabilidad del hardware. Al alcanzar la resolución de 112x84, la división en lotes evitó el error de \textit{Out of Memory}, consolidándose como una técnica obligatoria para problemas que exceden la memoria física.
    
    \item \textbf{Límites del \textit{Latency Hiding}:} La complejidad de coordinación y las dependencias de sincronización superaron al beneficio del solapamiento temporal. El enfoque asíncrono no mejoró el tiempo neto de ejecución, confirmando que su implementación debe justificarse por restricciones de memoria y no necesariamente por métricas de \textit{speedup} absoluto.
\end{itemize}

\end{document}

